\subsection{Influence of Retrieved Evidence}

We evaluate all $2^3$ combinations of API knowledge, repository context, and
similar code on ten execution-backed tasks for each backend. This yields 160
task--condition runs and 480 test-executed candidate generations. All source
conditions for a task share the same query decomposition and retrieval intent.
Table~\ref{tab:source_factorial_effects} reports task-paired marginal effects,
while Table~\ref{tab:source_interactions} reports difference-in-differences
interactions.

No individual source exhibits a statistically reliable marginal effect on
aggregate uncertainty after Holm correction. For GPT, the marginal uncertainty
effects of API, context, and similar-code evidence are
0.002,
0.003, and
-0.018,
respectively. The corresponding Gemini effects are
0.003,
0.009, and
-0.010.
The interaction analysis is more informative: context and similar code have a
positive GPT Pass@1 interaction of
33.3
percentage points ($p_{\mathrm{Holm}}=
0.032$),
whereas API and similar-code evidence reduce Gemini uncertainty jointly by
0.059
($p_{\mathrm{Holm}}=
0.032$).

\noindent\textbf{Finding 1.}
Retrieved evidence does not have a backend-independent additive effect on
uncertainty. Instead, uncertainty and correctness are shaped by interactions
among heterogeneous sources, supporting consensus-aware fusion rather than a
fixed universal source ordering.

\subsection{Uncertainty-Aware Mitigation}

We isolate query decomposition, uncertainty-aware filtering, guided generation,
verified candidate selection, and repair using six paired configurations per
task and backend. Table~\ref{tab:uncertainty_component_effectiveness}
distinguishes raw-sample performance from the final effective output.
OpenCoder increases effective Pass@1 from
10.0\% to
80.0\% for GPT and from
0.0\% to
60.0\% for Gemini. The paired gains are
70.0 points
(95\% CI [40.0, 100.0];
exact McNemar $p=0.016$) and
60.0 points
(95\% CI [30.0, 90.0];
exact McNemar $p=0.031$), with 7/0 and
6/0 improved/regressed task pairs.

The mitigation pathway differs by backend. For GPT, verified selection reaches
the same 80.0\% effective Pass@1 as the full
system, and repair does not recover an additional post-selection failure. For
Gemini, selection succeeds on 14.3\% of
initial failures and repair succeeds on 33.3\%
of attempted post-selection failures, raising final Pass@1 to
60.0\%. Calibration is not uniformly improved:
ECE changes from 0.162 to
0.222 for GPT, but from
0.191 to
0.153 for Gemini. Failure AUROC for full OpenCoder is
0.167 and 0.905, respectively.

\noindent\textbf{Finding 2.}
Uncertainty-aware decomposition, evidence control, verification, and repair
substantially improve final functional correctness, but raw uncertainty
calibration remains backend-dependent. OpenCoder is therefore supported as an
effective mitigation architecture; stronger cross-backend calibration remains
an explicit limitation rather than a claimed universal property.
